\chapter{Datenaufnahme}
\label{ch:datenaufnahme}

Die Qualität einer Behavior-Cloning-Policy ist durch ihre Daten begrenzt. Dieses Kapitel
beschreibt deshalb ausführlich, wie die Demonstrationen entstehen: vom Teachen der Punkte
über die Bahnplanung und das eingespeiste Rauschen bis zum fertigen Datensatz.

% ==============================================================================
\section{Grundentscheidung: Planer statt Teleoperation}
\label{sec:planer-statt-teleop}

\paragraph{Idee.} Üblich ist es, Demonstrationen von einem Menschen vorführen zu lassen, etwa
per Gamepad oder durch Handführung. Modul~4 trennt dagegen zwei Rollen: Der Mensch teacht nur
wenige Punkte der Bahn, die eigentliche Fahrt führt ein Bahnplaner aus, und dabei wird
aufgezeichnet.

\paragraph{Begründung.}
\begin{itemize}
  \item \textbf{Gleichmäßiges Tempo.} Die Policy lernt die Geschwindigkeit aus dem Abstand
    zweier Aktionen (\cref{sec:unumkehrbar}). Ein Mensch fährt jedes Mal anders schnell und
    zögert an schwierigen Stellen. Solche Pausen lernt die Policy als \enquote{stehen
    bleiben}. Der Planer fährt immer gleich.
  \item \textbf{Korrekturbeispiele.} Ein Planer fährt ohne Störung immer exakt die ideale
    Bahn, was die Fehlerfortpflanzung aus \cref{sec:covariate-shift} begünstigt. Weil er die
    ideale Bahn aber kennt, kann man ihn gezielt stören und trotzdem die richtige Aktion
    speichern (\cref{sec:noise-umsetzung}). Bei einem menschlichen Bediener ist die ideale
    Aktion unbekannt.
  \item \textbf{Aufwand.} Nach einmaligem Teachen nimmt die Anlage viele Episoden ohne
    Eingriff auf.
\end{itemize}

\paragraph{Preis.} Die Demonstrationen sind nur so gut wie die geteachten Punkte, und alle
Bahnen haben dieselbe Struktur. Die Vielfalt muss gezielt über Objektlagen, Rauschen und
Versatz hergestellt werden (\cref{sec:lagen}).

% ==============================================================================
\section{Ablauf und Punkte}
\label{sec:ablauf-punkte}

Der Ablauf einer Demonstration ist in einer Ablaufdatei als Folge von Punktnamen mit
Bewegungsart festgelegt. \Cref{fig:ablauf-punkte} zeigt die Folge, die für das Projekt
verwendet wird.

\begin{figure}[htbp]
  \centering
  \begin{tikzpicture}[font=\small, x=1cm, y=1cm,
      pt/.style={circle, fill=black, inner sep=1.4pt},
      lbl/.style={font=\scriptsize\ttfamily}]
    % Tisch, Objekt, Station
    \fill[black!10] (-0.5,0) rectangle (12.5,-0.25);
    \fill[blue!40] (2.6,0) rectangle (3.4,0.25);
    \fill[black!25] (10.2,0) rectangle (11.8,1.0);
    \node[font=\scriptsize] at (11,0.5) {Prüfvorr.};
    \node[font=\scriptsize, blue!50!black] at (3,-0.5) {Bauteil};
    % Punkte
    \node[pt] (start) at (-0.2,4.4) {}; \node[lbl, above] at (start) {CLEAR\_FOV};
    \node[pt] (ap1) at (3,3.3) {};    \node[lbl, left=1mm] at (ap1) {ANFLUG};
    \node[pt] (pre) at (3,1.8) {};    \node[lbl, left=1mm] at (pre) {PRE\_PICK};
    \node[pt] (pick) at (3,0.3) {};   \node[lbl, left=1mm] at (pick) {PICK};
    \node[pt] (place) at (11,2.2) {}; \node[lbl, above] at (place) {PRE\_PLACE};
    % Bahnen
    \draw[-{Stealth[length=2mm]}, thick] (start) to[out=-5, in=120] node[above right, font=\scriptsize] {PTP, mit Drehung} (ap1);
    \draw[-{Stealth[length=2mm]}, thick, red!70!black] ([xshift=-1mm]ap1.south) -- ([xshift=-1mm]pre.north);
    \draw[-{Stealth[length=2mm]}, thick, red!70!black] ([xshift=-1mm]pre.south) -- ([xshift=-1mm]pick.north);
    \draw[-{Stealth[length=2mm]}, thick, red!70!black] ([xshift=1mm]pick.north) -- ([xshift=1mm]pre.south);
    \node[font=\scriptsize, red!70!black, anchor=west] at (3.15,2.55) {LIN, \SI{5}{\centi\metre}};
    \node[font=\scriptsize, red!70!black, anchor=west] at (3.15,1.05) {LIN ab/auf, \SI{5}{\centi\metre}};
    \draw[-{Stealth[length=2mm]}, thick] (pre) to[out=50, in=170] node[above, font=\scriptsize] {PTP} (place);
    \node[font=\scriptsize, align=left, anchor=west] at (3.6,0.35) {schließen, \SI{0.5}{\second} warten};
    \draw[densely dotted] (10.2,2.2) -- (9.4,2.2);
    \draw[{Stealth[length=1.4mm]}-{Stealth[length=1.4mm]}] (9.6,1.0) -- node[left, font=\scriptsize] {\SI{7}{\centi\metre}} (9.6,2.2);
    \draw[densely dotted] (10.2,1.0) -- (9.4,1.0);
    \node[font=\scriptsize, green!40!black, anchor=west] at (11.9,2.2) {Übergabe an Modul~2};
  \end{tikzpicture}
  \caption[Ablauf einer Demonstration]{Ablauf einer Demonstration (Seitenansicht,
  schematisch). Die Startstellung \code{CLEAR_FOV} gibt das Bild der Draufsichtkamera
  frei. \code{ANFLUG} und \code{PRE_PICK} liegen \SI{15}{} bzw.\ \SI{5}{\centi\metre}
  senkrecht über \code{PICK} und werden daraus berechnet. Die gesamte Drehung geschieht in der
  Gelenkfahrt bis \code{ANFLUG}, darunter fährt der Arm linear und senkrecht.}
  \label{fig:ablauf-punkte}
\end{figure}

\paragraph{Abgeleitete Punkte.} Nur \code{CLEAR_FOV}, \code{PICK} und \code{PRE_PLACE} werden
geteacht. \code{ANFLUG} und \code{PRE_PICK} legt die Ablaufdatei relativ zu \code{PICK} fest.
Für eine neue Objektlage muss damit nur noch \code{PICK} bekannt sein, was das automatische
Umlegen erst möglich macht (\cref{sec:lagen}).

\paragraph{Punkte aus der Steuerung.} Die Ablaufdatei enthält nur Namen. Die Koordinaten
werden bei jedem Start frisch aus der Punktedatenbank der Steuerung gelesen. Wird ein Punkt
nachgeteacht, fährt schon die nächste Aufnahme die neue Lage an, ohne Export oder Neustart.
Ändert sich ein Punkt während einer laufenden Aufnahme, wird sie beendet, statt mit einem
veralteten Plan weiterzufahren.

\paragraph{Teachen.} Für das Teachen gibt es drei Wege, die nach Rückmeldungen aus der
Arbeit an der Anlage entstanden sind:
\begin{itemize}
  \item mit dem Gamepad über den Teleop-Port,
  \item per Handführung mit den Knöpfen am Arm: Der Punkt wird in der Ablauftabelle der
    Oberfläche markiert, der Arm frei geführt und der Punkt am Arm gespeichert (Teach-Serie),
  \item per Touch-up aus der Oberfläche: Die aktuelle Stellung des Arms wird mit einem Klick
    für den gewählten Punkt übernommen.
\end{itemize}
Außerdem kann die Startlage des Bauteils aus der Objekterkennung von Modul~1 übernommen
werden (\cref{sec:lagen}).

\paragraph{Tischpunkte.} Zusätzlich werden Punkte direkt auf dem Tisch geteacht, mittig über
den ArUco-Markern im $3 \times 3$-Raster. Aus ihnen wird die Tischoberfläche bestimmt
(\cref{sec:tisch}). Sie dient der Kollisionsprüfung, dem Ausrichten des Greifers senkrecht
zum Tisch und dem automatischen Umlegen des Bauteils.

% ==============================================================================
\section{Bahnplanung}
\label{sec:bahnplanung}

Der Planer erzeugt aus den Punkten eine Bahn, die mit dem Takt von \SI{15}{\hertz} abgetastet
ist. Die Geschwindigkeit steckt damit im Abstand zweier aufeinanderfolgender Sollwerte.

\paragraph{Geschwindigkeitsprofil.} Jeder Abschnitt beginnt mit einer Anfahrrampe von
\SI{0.5}{\second}, fährt mit konstanter Geschwindigkeit und bremst bis zum Stillstand. In
einer ersten Fassung sprang die Bahn direkt auf volle Geschwindigkeit; der Regler der
Steuerung schwang daraufhin in der \ac{vm} um bis zu \SI{0.1}{rad} über. Die Rampen sind fest
und in jeder Episode gleich, das gelernte Tempo bleibt also konsistent.
\Cref{tab:geschwindigkeiten} nennt die Geschwindigkeiten.

\begin{table}[htbp]
  \centering
  \caption{Geschwindigkeiten des Bahnplaners}
  \label{tab:geschwindigkeiten}
  \begin{tabular}{@{}l r@{}}
    \toprule
    Abschnitt & Geschwindigkeit \\
    \midrule
    Transfer, linear & \SI{0.15}{\metre\per\second} \\
    Gelenkfahrt (\acs{ptp}) & \SI{0.5}{rad/s} je Gelenk \\
    Endanflug zum Greifpunkt & \SI{0.02}{\metre\per\second} \\
    Override der Steuerung & \num{0.5} \\
    \bottomrule
  \end{tabular}
\end{table}

Der Endanflug wurde nach dem Anlagentest vom 16.11.\ von \SI{0.05}{} auf
\SI{0.02}{\metre\per\second} gesenkt. Langsameres Greifen gibt der Policy mehr Takte mit
Bildern unmittelbar vor dem Bauteil und verringert den Versatz des Bauteils im Greifer.

\paragraph{Überschleifen.} An Durchfahrpunkten hält der Arm nicht an. Das Folgesegment
beginnt bereits während der Bremsrampe des vorherigen, Brems- und Anfahrrampe überlagern
sich, und die Geschwindigkeit bleibt stetig. Zwei Gelenkfahrten werden im Gelenkraum
überlagert, jede Kombination mit einer Linearfahrt kartesisch. Nie überschliffen wird am
Start, am Ende und an Greiferwechseln. Überschliffen wird an \code{ANFLUG} (\SI{5}{\centi\metre}) und an \code{PRE_PICK}
(\SI{1}{\centi\metre} hinab, \SI{4}{\centi\metre} hinauf). Das Überschleifen an der
Vorgreifposition war eine Folge
der ersten Versuche in der \ac{vm}: Dort endete eine Gelenkfahrt im Stillstand, und die Policy
blieb genau dort stehen, weil \enquote{angekommen} und \enquote{gleich losfahren} im Zustand
nicht zu unterscheiden waren (\cref{sec:multimodal}).

\paragraph{Greifer.} Am Greifpunkt hält der Arm, der Greifer schließt, und die Bahn wartet
\SI{0.5}{\second}, bevor sie anhebt. Die Wartezeit sorgt dafür, dass der Arm erst anhebt, wenn
die Backen geschlossen sind. Da der Greifer keine Rückmeldung gibt, ist der Greiferzustand in
den Daten der \emph{kommandierte} Wert.

\todo[inline]{Greifer-Totzeit (\SI{0.5}{\second}) an der Anlage per Video bestätigen.}

% ==============================================================================
\section{Eingespeistes Rauschen}
\label{sec:noise-umsetzung}

\subsection{Idee}

Die Grundlage ist \ac{dart} mit asymmetrischen Datenpaaren (\cref{sec:rauschen-mathe}):
Während der Fahrt wird die Bahn ausgelenkt, gespeichert werden der ausgelenkte Zustand mit
seinen Bildern und als Aktion der nächste Punkt der idealen Bahn. Damit die Zuordnung später
überprüft und bei Bedarf neu gebildet werden kann, speichert der Datensatz zusätzlich ideale
und tatsächliche Bahn getrennt.

\subsection{Kartesisch statt im Gelenkraum}

Das Rauschen wird kartesisch auf die \ac{tcp}-Pose gegeben und anschließend mit der inversen
Kinematik in Gelenkwinkel umgerechnet. Einfacher wäre Rauschen direkt auf die Gelenkwinkel,
aber dann hängt die Auslenkung des Werkzeugs über die Jacobi-Matrix von der Armstellung ab:
Dasselbe Rauschen auf einem schulternahen Gelenk verschiebt das Werkzeug viel stärker als auf
einem Handgelenk. Eine Grenze in einer festen Richtung, etwa \enquote{höchstens wenige
Millimeter zum Tisch hin}, ließe sich so nicht garantieren. Die kartesische Variante ist
aufwändiger, aber die einzige, deren Amplitude und Richtung kontrollierbar sind.

\subsection{Amplitude und zeitlicher Verlauf}

Position und Drehung werden mit je einem geglätteten \ac{ou}-Prozess
(\cref{eq:ou-diskret}) ausgelenkt. \Cref{tab:rauschen} nennt die Parameter.

\begin{table}[htbp]
  \centering
  \caption{Parameter des eingespeisten Rauschens}
  \label{tab:rauschen}
  \begin{tabularx}{\textwidth}{@{}l r X@{}}
    \toprule
    Parameter & Wert & Begründung \\
    \midrule
    Amplitude Position & \SI{15}{\milli\metre} & im Bild deutlich sichtbar, über Kamera-
      und Schwingungsrauschen \\
    Amplitude Drehung & \SI{0.05}{rad} & entsprechend \\
    Zeitkonstante $\tau$ & \SI{0.8}{\second} & langsames Schlingern statt Zittern \\
    Glättung (2.\,Ordnung) & \SI{0.25}{\second} & begrenzte Geschwindigkeit und
      Beschleunigung der Störung \\
    Stärke je Episode & gleichverteilt $0 \dots 1$ & ruhige und stark gestörte Fahrten
      im selben Datensatz \\
    \bottomrule
  \end{tabularx}
\end{table}

Eine Auslenkung im Millimeterbereich wäre im Bild von Bildrauschen kaum zu unterscheiden und
würde kein lernbares Korrektursignal erzeugen. Deshalb ist die Amplitude bewusst im
Zentimeterbereich gewählt. Der ungeglättete \ac{ou}-Prozess hat allerdings keine begrenzte
Geschwindigkeit; die Bahnen hätten einzelne Gelenke kurzzeitig über ihre zulässige
Geschwindigkeit gebracht und wären verworfen worden. Ein nachgeschaltetes, kritisch
gedämpftes Glied zweiter Ordnung behält den Charakter der Störung bei, macht aber
Geschwindigkeit und Beschleunigung endlich. Die Verstärkung ist so gewählt, dass die
Standardabweichung trotzdem der Amplitude entspricht.

Die Rauschstärke wird je Episode zufällig gezogen. Eine Ersatzstudie ohne Bilder ergab, dass
diese Variante bei seitlichen Stößen am robustesten ist und zugleich die ruhigsten
Bewegungen erzeugt (\cref{ch:versuche}).

\subsection{Trichter vor dem Greifpunkt}

Kurz vor dem Greifen muss die Bahn genau sein, denn das Spiel zwischen Bauteil und Backen
beträgt nur etwa \SI{2.5}{\milli\metre} je Seite (Bauteil \SI{39.7}{\milli\metre}, Backen
\SI{44.8}{\milli\metre}). Das Rauschen wird deshalb mit dem Abstand $d$ zum Greifpunkt
gedämpft:
\begin{equation}
  \delta_t' = s(d_t)\,\delta_t,\qquad
  s(d) = \begin{cases}
    0, & d \le d_0,\\
    h\!\left(\dfrac{d - d_0}{d_1 - d_0}\right), & d_0 < d < d_1,\\
    1, & d \ge d_1,
  \end{cases}
  \label{eq:trichter}
\end{equation}
mit $d_0 = \SI{2}{\centi\metre}$, $d_1 = \SI{15}{\centi\metre}$ und einer glatten Rampe $h$
(\cref{fig:trichter}). Nach unten, zum Tisch hin, ist die Auslenkung zusätzlich begrenzt.

\begin{figure}[htbp]
  \centering
  \begin{tikzpicture}[font=\small, x=1cm, y=1cm]
    \fill[blue!8] (0,1) -- (6.5,1) -- (9.3,0) -- (6.5,-1) -- (0,-1) -- cycle;
    \draw[blue!40, densely dashed] (0,1) -- (6.5,1) -- (9.3,0) -- (6.5,-1) -- (0,-1);
    \draw[thick] (0,0) -- (10,0);
    \fill (0,0) circle (2pt) node[left] {Start};
    \fill (10,0) circle (2pt) node[right] {Greifpunkt};
    \draw[red!70!black, thick] plot[smooth, tension=0.7] coordinates
      {(0,0) (0.9,0.45) (1.8,0.75) (2.7,0.3) (3.6,-0.5) (4.5,-0.8) (5.4,-0.2) (6.3,0.55)
       (7.2,0.4) (8.1,-0.15) (8.8,0.05) (9.4,0) (10,0)};
    \foreach \x/\y in {1.8/0.75, 4.5/-0.8, 6.3/0.55} {
      \draw[-{Stealth[length=1.6mm]}, green!50!black, thick] (\x,\y) -- (\x+0.7,0.04*\y);
    }
    \draw[{Stealth[length=1.4mm]}-{Stealth[length=1.4mm]}] (0.3,-1) -- (0.3,1);
    \node[font=\scriptsize, anchor=west] at (0.35,1.2) {Amplitude};
    \draw[densely dotted] (6.5,-1.1) -- (6.5,-1.5);
    \draw[densely dotted] (9.3,-0.1) -- (9.3,-1.5);
    \draw[densely dotted] (10,-0.1) -- (10,-1.5);
    \draw[{Stealth[length=1.4mm]}-{Stealth[length=1.4mm]}] (6.5,-1.4) -- node[below, font=\scriptsize] {$d_1 = \SI{15}{\centi\metre}$} (10,-1.4);
    \node[font=\scriptsize, anchor=north] at (9.65,-1.6) {$d_0$};
    \begin{scope}[shift={(0,-2.5)}, font=\scriptsize]
      \draw[thick] (0,0) -- (0.6,0); \node[anchor=west] at (0.7,0) {ideale Bahn (Aktion)};
      \draw[red!70!black, thick] (4.0,0) -- (4.6,0); \node[anchor=west] at (4.7,0) {gefahrene Bahn (Beobachtung)};
      \draw[-{Stealth[length=1.6mm]}, green!50!black, thick] (9.3,0) -- (9.9,0); \node[anchor=west] at (10.0,0) {Korrektur};
    \end{scope}
  \end{tikzpicture}
  \caption[Rauschkorridor um die ideale Bahn]{Rauschkorridor um die ideale Bahn
  (schematisch). Bis $d_1$ vor dem Greifpunkt ist die Auslenkung voll wirksam, danach wird sie
  gedämpft und ist ab $d_0$ null. Aus der Paarung von ausgelenktem Zustand und idealer
  Sollstellung lernt die Policy, zur Bahn zurückzusteuern.}
  \label{fig:trichter}
\end{figure}

\subsection{Problem: blinder Endanflug und gezielter Versatz}
\label{sec:versatz}

\paragraph{Problem.} Der Trichter hat eine unerwünschte Folge: Im senkrechten Abstieg auf das
Bauteil ist das Rauschen praktisch null (ab etwa \SI{5}{\centi\metre} unter
\SI{2}{\milli\metre}, ab \SI{3}{\centi\metre} unter \SI{0.3}{\milli\metre}). Die Policy sieht
in den Daten also nie, was zu tun ist, wenn sie versetzt oder verdreht über dem Bauteil
ankommt. Genau das wurde an der Anlage beobachtet: Die Policy stand \SIrange{2}{16}{\milli\metre}
und bis etwa \ang{50} neben dem Bauteil und korrigierte im Anflug kaum noch.

\paragraph{Lösung.} Ein Teil der Episoden fährt mit einem gezielten seitlichen Versatz und
einer Drehung über dem Greifpunkt an, die im Abstieg bis \SI{2}{\centi\metre} über dem
Greifpunkt auf null abklingen (\cref{fig:versatz}). Die Aktion bleibt die ideale Bahn und
zeigt damit zur Mitte zurück. Die aktuellen Werte sind in \cref{tab:versatz} aufgeführt.

\begin{table}[htbp]
  \centering
  \caption[Gezielter Versatz im Anflug]{Gezielter Versatz im Anflug (Stand 19.11.2026)}
  \label{tab:versatz}
  \begin{tabular}{@{}l r@{}}
    \toprule
    Parameter & Wert \\
    \midrule
    Anteil der Episoden mit Versatz & \SI{50}{\percent} \\
    seitlich, Richtung gleichverteilt & \SIrange{5}{10}{\milli\metre} \\
    Drehung um die Senkrechte & \SIrange{10}{20}{\degree} \\
    Ausreißer (je \SI{15}{\percent}) & bis \SI{25}{\milli\metre} bzw.\ \ang{35} \\
    Auf- und Abbau höchstens & \SI{30}{\milli\metre\per\second}, \SI{12}{\degree\per\second} \\
    voll wirksam & \SIrange{5}{10}{\centi\metre} vor dem Greifpunkt \\
    null ab & \SI{2}{\centi\metre} vor dem Greifpunkt \\
    \bottomrule
  \end{tabular}
\end{table}

\paragraph{Folgeproblem.} In einer ersten Fassung wurde die Drehung zu schnell auf- und
wieder abgebaut. Gelenk~6 kam der Umkehr nicht nach, lief bis zu \SI{0.34}{rad} hinterher, und
die Episoden wurden verworfen. Seitdem ist die Änderungsrate begrenzt, der Verlauf zusätzlich
über \SI{0.8}{\second} geglättet, und die Drehung beginnt früher (\SI{30}{\centi\metre} vor
dem Greifpunkt), damit sie eine Haltephase hat.

\begin{figure}[htbp]
  \centering
  \begin{tikzpicture}[font=\small, x=0.75cm, y=2.4cm]
    \draw[-{Stealth[length=2mm]}] (0,0) -- (16.3,0);
    \node[font=\scriptsize, anchor=north east] at (16.3,-0.15) {Abstand zum Greifpunkt};
    \draw[-{Stealth[length=2mm]}] (0,0) -- (0,1.25) node[above] {Versatz / Rauschen (relativ)};
    \foreach \x/\l in {0/0, 2/2, 5/5, 10/10, 15/15} {
      \draw (\x,0) -- (\x,-0.04) node[below, font=\scriptsize] {\l\,cm};
    }
    \draw[blue!60!black, thick] (15.5,1) -- (15,1) .. controls (11,1) and (5,0.2) .. (2,0) -- (0,0);
    \node[blue!60!black, font=\scriptsize, anchor=south east] at (15.5,1.02) {Trichter des Rauschens};
    \draw[orange!80!black, thick] (15.5,0) -- (15,0) .. controls (13,0) and (12,1) .. (10,1)
      -- (5,1) .. controls (3.5,1) and (3,0) .. (2,0) -- (0,0);
    \node[orange!80!black, font=\scriptsize, anchor=south] at (7.5,1.02) {gezielter Versatz};
    \fill[red!10] (0,0) rectangle (2,1.2);
    \node[font=\scriptsize, align=center] at (1,0.6) {Backen\\am\\Bauteil};
  \end{tikzpicture}
  \caption[Trichter und gezielter Versatz über dem Abstand zum Greifpunkt]{Trichter des
  Rauschens (blau) und gezielter Versatz (orange) über dem Abstand zum Greifpunkt
  (schematisch). Der Versatz schließt die Lücke, die der Trichter im Endanflug hinterlässt. In
  den letzten \SI{2}{\centi\metre} sind beide null.}
  \label{fig:versatz}
\end{figure}

% ==============================================================================
\section{Kollisionsprüfung und Neuziehen}
\label{sec:kollision}

Jede geplante Bahn wird vor der Fahrt vollständig geprüft. Verletzt sie eine Bedingung, wird
das Rauschen neu gezogen, höchstens 20-mal (Rejection Sampling). Geprüft werden:
\begin{itemize}
  \item jede Stellung gegen die Achsgrenzen und gegen eine Lösung der inversen Kinematik,
  \item die Änderung der Gelenkwinkel zwischen zwei Takten, um Sprünge nahe singulärer
    Stellungen zu erkennen, auch wenn die inverse Kinematik formal eine Lösung liefert,
  \item der Abstand von Werkzeug, Kamera, Handgelenk und Ellbogen zu Tisch und Hindernissen
    (\cref{sec:tisch}) und
  \item die seitliche Abweichung in den letzten \SI{2}{\centi\metre} vor dem Greifen
    (höchstens \SI{1.5}{\milli\metre}, \cref{sec:versatz}).
\end{itemize}
Die inverse Kinematik läuft dabei offline bei der Planung und nicht im Echtzeittakt. Die
Prüfung braucht ohnehin die ganze Bahn, und der Roundtrip zur Steuerung geht so nicht in die
Taktzeit ein.

\subsection{Tisch und Hindernisse}
\label{sec:tisch}

\paragraph{Idee.} Das erste Kollisionsmodell bestand aus Kugeln um Werkzeug, Handgelenk und
Ellbogen und einem pauschalen Tisch \SI{15}{\centi\metre} unter dem tiefsten Punkt.

\paragraph{Probleme.} Der pauschale Tisch fing nichts ab. Und die Kugeln waren für diese Zelle
zu grob: Der Übergabepunkt liegt \SI{7}{\centi\metre} über der Prüfvorrichtung, und Tischpunkte
an der Kamerastange liegen nur wenige Zentimeter davor. Mit Kugeln wären beide eine
\enquote{Kollision}.

\paragraph{Lösung.} Die Tischoberfläche wird aus den geteachten Tischpunkten bestimmt: ab
drei Punkten als Ebene $z = ax + by + c$, ab genügend Punkten zusätzlich mit einem
quadratischen Anteil für den leicht gewölbten Tisch. Der quadratische Anteil wird nur
verwendet, wenn er die Punkte deutlich besser trifft als die Ebene; sonst würde er Streuung
beim Teachen nachbilden. Die Restabweichung an den Punkten (bis \SI{0.4}{\milli\metre}) wird
zwischen den Punkten interpoliert und addiert, weil das Bauteil beim Umlegen über
\SIrange{10}{30}{\centi\metre} sonst spürbar zu hoch oder zu tief abgelegt wurde. Feste
Hindernisse (Kamerastange, Prüfvorrichtung, Roboterbasis) sind als Quader beschrieben, jeweils
mit eigenem Mindestabstand je Stützpunkt des Arms \quelle{einstellungen/hindernisse.json}.
Sie gelten nur an der realen Anlage und nur für Modul~4.

% ==============================================================================
\section{Kameras und Bildaufnahme}
\label{sec:kameras}

\subsection{Kamerawahl}

\paragraph{Handgelenk und Draufsicht.} Eine Policy nur mit Handgelenkkamera wäre unabhängig
von der Lage einer festen Kamera, ihr fehlt aber der Überblick: Liegt das Bauteil nicht im
Bild der Startstellung, findet sie es nicht, und kurz vor dem Greifen verdecken die Backen
einen großen Teil des Bildes. Die Draufsicht liefert den Überblick. Aufgezeichnet werden
deshalb beide Kameras. Ein Vergleich nur mit Handgelenkkamera bleibt jederzeit möglich, weil
sich ein vorhandener Bildstrom beim Training weglassen, ein fehlender aber nicht nachträglich
ergänzen lässt.

\paragraph{Objektiv.} Zur Wahl standen für die Handgelenkkamera ein rektilineares Objektiv mit
festem Fokus und ein Fischauge mit \SI{1.85}{\milli\metre} Brennweite. Das Festfokus-Objektiv
wäre nur im Endanflug scharf, also gerade dort, wo die Backen das Bild verdecken, und im
Anflug unscharf, wo das Korrektursignal aus dem Rauschen seinen Wert hat. Entschieden wurde
für das Fischauge: Unschärfe ist unwiederbringlicher Informationsverlust, die Verzeichnung
dagegen in jedem Bild gleich und wird vom Netz mitgelernt. Der Preis ist eine geringere
effektive Auflösung; das Bauteil belegt im $240 \times 320$-Bild nur wenige Pixel.

\paragraph{Anbindung.} Die Daheng-Kameras sind \acs{usb}3-Vision-Kameras und lassen sich nicht
über die übliche Videoschnittstelle öffnen, sondern nur über das Galaxy SDK des Herstellers.
Die Python-Schnittstelle des SDK ist kein installierbares Paket; der Adapter findet sie zur
Laufzeit. Weitere Stolpersteine der Inbetriebnahme sind im Dokument
\code{Dokumentation/Kamera-Inbetriebnahme.md} beschrieben, unter anderem ein in der Kamera
gespeicherter Bildausschnitt, der ohne aktives Setzen stumm eine Bildecke lieferte.

\subsection{Einstellungen}

Die Policy lernt die Bildstatistik mit. Jede Automatik, die Helligkeit oder Farbe zwischen
Sitzungen verändert, verschiebt daher die Eingangsverteilung. Die Einstellungen wurden
deshalb schrittweise festgelegt:
\begin{description}
  \item[Bildrate] beide Kameras \SI{60}{fps}, direkt an getrennten \acs{usb}-3-Anschlüssen.
    An einem Verteiler erreichten sie nur \SI{23}{fps}, mit \SI{30}{fps} war das neueste Bild
    bis zu \SI{33}{\milli\second} alt und lag damit über dem Budget aus \cref{sec:sync}.
  \item[Belichtung und Verstärkung] fest, einstellbar im Fenster \enquote{Kameraeinstellungen}
    und je Aufnahme protokolliert.
  \item[Weißabgleich] fest je Kamera (Seriennummer). Die Kamera kennt nur \enquote{aus} und
    \enquote{einmal messen}; beim einmaligen Messen sprang der Blauanteil zwischen Sitzungen
    von \num{1.85} auf \num{2.68}.
  \item[Kameralicht] Seit dem 17.11.\ ist eine Leuchte an der Handgelenkkamera Standard. Das
    Bild beim Greifen war vorher oft zu dunkel; mit Licht ist es etwa dreimal heller und
    kontrastreicher.
  \item[Zuordnung] Bei zwei Daheng-Kameras ist die Reihenfolge der Geräte nicht stabil. Beide
    werden deshalb über ihre Seriennummer zugeordnet, sonst könnten Handgelenk- und
    Draufsichtbild unbemerkt vertauscht im Datensatz landen.
\end{description}

\subsection{Draufsicht ohne Rektifizierung}

\paragraph{Problem.} Die Policy lernt \enquote{Bauteil an dieser Bildposition, also diese
Bewegung}. Verschiebt sich die Draufsichtkamera zwischen Aufnahme und Ausführung, greift sie
systematisch daneben, ohne dass ein Fehler gemeldet wird. Die ArUco-Marker auf dem Tisch
helfen dabei nicht von selbst, denn die Policy nutzt sie nicht als Bezug.

\paragraph{Erste Lösung, verworfen.} Ursprünglich sollte jedes Bild über eine Homographie aus
den Markern auf eine feste Draufsicht entzerrt werden. Für alles in der Tischebene wäre die
Kameraposition dann weitgehend unerheblich. Die Prüfvorrichtung, das Ziel der Fahrt, liegt
aber außerhalb der Marker und ist \SI{70}{\milli\metre} hoch. Eine Homographie der Tischebene
hätte genau diesen Bildteil verzerrt oder abgeschnitten. Die Entscheidung wurde am 11.11.\
revidiert: Die Bilder werden unverändert aufgezeichnet und verwendet.

\paragraph{Stattdessen.} Gegen eine verschobene Kamera wirkt die bewusste Variation der
Kameraposition im Datensatz; jede Episode trägt die Kameraposition als Metadatum. Ein
Referenzbild erlaubt zu prüfen, ob die Kamera weiter verschoben ist als in den Daten
abgedeckt. Die Empfindlichkeit der Policy gegenüber der Kameraposition wurde gesondert
untersucht (\cref{ch:versuche}).

\subsection{Synchronisation}
\label{sec:sync}

Kameras und Roboterzustand liefern asynchron mit verschiedenen Raten. Jede Quelle läuft in
einem eigenen Thread und hält ihren jüngsten Wert mit Zeitstempel bereit. Im
\SI{15}{\hertz}-Takt greift der Recorder von jeder Quelle den jüngsten Wert ab und vergleicht
die Zeitstempel. Liegen Bild und Zustand mehr als \SI{30}{\milli\second} auseinander, also
weniger als einen halben Takt, wird der Takt als ungültig markiert. Sind mehr als
\SI{5}{\percent} der Takte einer Episode markiert, wird die Episode verworfen.

% ==============================================================================
\section{Objektlagen}
\label{sec:lagen}

Damit die Policy das Bauteil an beliebigen Stellen findet, braucht der Datensatz viele
verschiedene Lagen. Für jede Lage den Greifpunkt von Hand nachzuteachen, wäre der größte
Zeitaufwand der Aufnahme. Es wurde in zwei Schritten gelöst.

\paragraph{Zurücklegen.} Nach jeder Episode legt der Arm das Bauteil am Greifpunkt wieder
ab und fährt zum Start, ohne Aufzeichnung und ohne Rauschen. Die Rückfahrt ist die Bahn
rückwärts, es wird also nur Geometrie gefahren, die schon geplant und geprüft war. So lässt
sich eine Lage mehrfach mit neuem Rauschen aufnehmen. Eine Abschätzung auf vorhandenen Daten
zeigte, dass der Abstand einer neuen Abweichung zum nächsten Trainingszustand mit der Zahl
$K$ der Wiederholungen auf \SI{73}{\percent} ($K=2$), \SI{60}{\percent} ($K=3$) und
\SI{53}{\percent} ($K=4$) sinkt und danach kaum noch. Wiederholungen ersetzen aber keine neue
Lage, weil das Bauteil im Bild an derselben Stelle bleibt.

\paragraph{Automatisches Umlegen.} Deshalb legt der Arm das Bauteil nach den Fahrten einer
Lage an eine \emph{neue} Lage. Die neue Position liegt innerhalb der um \SI{4}{\centi\metre}
geschrumpften Hülle der Tischpunkte, außerhalb von Sperrzonen und mindestens
\SI{3}{\centi\metre} von allen bisherigen Lagen entfernt. Gezogen wird über eine
Abdeckungskarte, weil rein zufälliges Ziehen manche Stellen häufte. Die Drehung ist
gleichverteilt in \ang{\pm 45}. Da der Roboter die Lage selbst hergestellt hat, ist sie exakt
bekannt und wird als Zahl in den Metadaten der Episode gespeichert.

\paragraph{Problem Drehbias.} Die Drehung wurde zunächst um die geteachte Lage verteilt.
Gemessen gegen die Startstellung der Policy waren die Demonstrationen dadurch einseitig
(\SI{62}{\percent} in die eine, \SI{25}{\percent} in die andere Richtung), und die Policy drehte
an der Anlage systematisch zu weit (\cref{ch:versuche}). Seit dem 17.11.\ wird die Drehung um
die Startstellung verteilt.

\paragraph{Start mit Modul~1.} Für die erste Lage einer Aufnahme kann die Lage aus der
Objekterkennung von Modul~1 übernommen werden. Der erste Griff wird dann nicht aufgezeichnet;
danach greift der Arm nur noch dort, wo er selbst abgelegt hat.

\paragraph{Höhenprotokoll.} An jedem Greiferwechsel werden Soll- und Ist-Höhe, die Tischhöhe
an der Stelle und die Bewegung des Arms beim Schaltbefehl protokolliert. Anlass war ein
Ablegen, das nach einigen Episoden zu tief geriet, obwohl alle Sollhöhen gleich waren.

% ==============================================================================
\section{Gültigkeit einer Episode}
\label{sec:gueltigkeit}

Eine Episode wird verworfen, wenn
\begin{itemize}
  \item ein Schutzstopp der Steuerung auslöst,
  \item der Arm vier Takte lang mehr als \SI{0.30}{rad} hinter der Sollstellung liegt
    (Schleppfehler),
  \item zu viele Takte die Synchronisation verletzen (\cref{sec:sync}) oder
  \item der Bediener sie verwirft.
\end{itemize}
Verworfene Episoden werden nicht gespeichert; die Aufnahme setzt mit der nächsten fort.

% ==============================================================================
\section{Datensatzformat und Export}
\label{sec:datensatz}

\paragraph{Zwei Formate.} Aufgezeichnet wird in einem eigenen, neutralen Format (NumPy-Archiv
und JSON je Episode), unabhängig von LeRobot und ohne Videokodierung im Takt. Die
Kodierung der Bilder als Video würde den Rechner während der Aufnahme belasten, und ein
Wechsel der LeRobot-Version würde sonst bereits aufgenommene Daten betreffen. Für das
Training werden die Episoden anschließend in das LeRobot-Format (Version~0.6.1, Format~3.0)
exportiert.

\paragraph{Inhalt.} \Cref{tab:schema} zeigt die Felder. Neben Zustand, Aktion und Bildern
werden ideale und tatsächliche Bahn getrennt gespeichert, dazu ein Gültigkeitsflag je Takt.

\begin{table}[htbp]
  \centering
  \caption{Felder einer Episode}
  \label{tab:schema}
  \small
  \begin{tabularx}{\textwidth}{@{}l l X@{}}
    \toprule
    Feld & Größe & Inhalt \\
    \midrule
    \code{observation.state} & 14 & sechs Gelenkwinkel (rad), \acs{tcp}-Position (m),
      \acs{tcp}-Quaternion, Greifer (0/1) \\
    \code{observation.images.wrist} & $240 \times 320 \times 3$ & Handgelenkkamera \\
    \code{observation.images.scene} & $240 \times 320 \times 3$ & Draufsichtkamera \\
    \code{action} & 7 & Gelenkwinkel der idealen Bahn im nächsten Takt, Greifer \\
    \code{aux.joints_ideal}, \code{aux.pose_ideal} & 6, 7 & ideale Bahn \\
    \code{aux.pose_noisy} & 7 & tatsächlich geplante, verrauschte Bahn \\
    \code{aux.sync_ok} & 1 & Takt innerhalb des Zeitbudgets \\
    \bottomrule
  \end{tabularx}
\end{table}

\paragraph{Metadaten.} Jede Episode speichert die Einstellungen ihrer Aufnahme: Punkte,
Objektlage, Rauschstärke und -parameter, Versatz, Geschwindigkeiten, Override,
Kameraeinstellungen und Kameraposition. Der Export schreibt Schemaversion, LeRobot-Version und
Quellordner in eine eigene Datei. Training und Ausführung prüfen dagegen und brechen bei
Abweichungen ab (\cref{sec:unumkehrbar}).
